\documentclass[../main.tex]{subfiles}

\begin{document}
\chapter{Riemannian Submanifolds}

We apply the theory of curvature to Riemannian submanifolds, and then use these concepts to derive a precise quantitative interpretation of the curvature tensor.

\section{The Second Fundamental Form}
Suppose $(M,g)$ is a Riemannian submanifold of a Riemannian manifold $(\tilde{M},\tilde{g})$, so that $g = \iota_M^*\tilde{g}$, where $\iota_M:M\hookrightarrow \tilde{M}$ is the inclusion map. Although we focus our attention in this chapter on embedded submanifolds for simplicity, the results we present are all local, so they apply in much greater generality.

\begin{remark}
	For example, if $M$ is an immersed submanifold of $\tilde{M}$, then for each $p\in M$ there exists an open neighborhood $U\subset M$ of $p$ such that $\iota_M|_U:U\to \tilde{M}$ is an embedding. In this case, we can apply the results of this chapter to the embedded submanifold $\iota_M(U)\subset \tilde{M}$.
\end{remark}

NOTICE: The results in the first section of this chapter apply virtually without modification to Riemannian submanifolds of pseudo-Riemannian manifolds (ones on which the induced metric is positive definite), so we state most of our theorems in that case. Also, most of these results can be adapted to manifolds and submanifolds with boundary, simply by embedding everything in slightly larger manifolds without boundary, but one might need to be careful about the statements of some of the results when the submanifold intersects the boundary. Since the interaction of submanifolds with boundaries is not our primary concern here, for simplicity we state all of these results in the case of empty boundary only.

Throughout this chapter, therefore, we assume that $(\tilde{M},\tilde{g})$ is a pseudo-Riemannian manifold of dimension $m$, and $(M,g)$ is an embedded $n$-dimensional Riemannian submanifold of $(\tilde{M},\tilde{g})$.

\end{document}
